\subsection{PRINCIPAL PARTS AND ASYMPTOTIC EXPANSIONS}

Let E be a comparison scale formed by functions with values in a non-discrete valued field K. Let V be a normed space over K, and let f be a function in \(\mathcal{H}(\mathfrak{F}, V)\) ; if there are a function \(g \in E\) and an element \(a \neq 0\) in V such that \(f \sim ag\) , one says that ag is a principal part of f relative to the scale E. From def. 1 of V, p.~10, f can have only one principal part relative to E, for if \(g_{1}\) , \(g_{2}\) are two functions in E and \(a_{1}\) , \(a_{2}\) are two \(\neq 0\) elements of V, then the relation \(a_{1}g_{1} \sim a_{2}g_{2}\) implies \(|g_{1}| \asymp |g_{2}|\) and consequently \(g_{1} = g_{2}\) , whence \((\mathbf{a}_{2} - \mathbf{a}_{1})g_{1} \ll g_{1}\) , and since \(g_{1}\) is not identically zero on any set in E this implies \(a_{2} = a_{1}\) .

If f has a principal part relative to a comparison scale E it has the same principal part relative to every comparison scale \(E' \supset E\) .

Examples. 1) For x real (resp. complex) tending to \(+\infty\) (resp. \(\infty\) ), every polynomial \(a_{0}x^{n} + a_{1}x^{n-1} + \ldots + a_{n}\) with coefficients in V, such that \(a_{0} \neq 0\) , has principal part \(a_{0}x^{n}\) with respect to the scale \(x^{n}\) (or any scale containing the \(x^{n}\) ). It follows that every rational fraction \(\frac{a_{0}x^{n} + \cdots + a_{m}}{b_{0}x^{n} + \cdots + b_{n}}\) with real or complex coefficients such that \(a_{0}b_{0} \neq 0\) has principal part \(\frac{a_{0}}{b_{0}}x^{m-n}\) with respect to the same scale.

\begin{enumerate}
\def\labelenumi{\arabic{enumi})}
\setcounter{enumi}{1}
\tightlist
\item
  A function may be comparable to all the functions of a scale and yet not admit a principal part with respect to this scale. For example, for \(x\) real tending to \(+\infty\) , \(\sqrt{x}\) has no principal part with respect to the scale \(x^n\) where \(n\) is a rational integer; \(\log x\) has no principal part with respect to the scale of \(x^\alpha\) ( \(\alpha\) arbitrary real); \(\exp(\sqrt{\log x})\) and \(x^x = e^{x\log x}\) have no principal part with respect to the scale of the \(x^\alpha (\log x)^\beta\) , nor with respect to the scale of the \(\exp(p(x))\) ( \(p\) a polynomial with no constant term).
\end{enumerate}

The concept of principal part admits extensive generalization. Suppose that a function \(\mathbf{f} \in \mathcal{H}(\mathfrak{F}, V)\) has a principal part \(\mathbf{a}_1g_1\) with respect to a scale \(\mathcal{E}\) ; the relation \(\mathbf{f} \sim \mathbf{a}_1g_1\) is equivalent to \(\mathbf{f} - \mathbf{a}_1g_1 \ll g_1\) (V, p.~216, def. 4); to study the function \(\mathbf{f}\) more closely one is thus led to consider the function \(\mathbf{f} - \mathbf{a}_1g_1\) . If this function has a principal part \(\mathbf{a}_2g_2\) with respect to \(\mathcal{E}\) one must have \(g_2 \ll g_1\) and \(\mathbf{f} - \mathbf{a}_1g_1 - \mathbf{a}_2g_2 \ll g_2\) .

More generally, suppose that the scale E is written parametrically in the form \((g_{\alpha})\) where \(\alpha\) runs through a set of indices A endowed with a totally ordered structure isomorphic to the opposite of the order structure of E: the relation \(\alpha < \beta\) is thus equivalent to \(g_{\beta} \ll g_{\alpha}\) . In these circumstances:

DEFINITION 2. One says that a function \(\mathbf{f} \in \mathcal{H}(\mathfrak{F}, V)\) admits an asymptotic expansion to precision \(g_{\alpha}\) (relative to the scale \(\mathcal{E}\) ) if there exists a family \((\mathbf{a}_{\lambda})_{\lambda \leqslant \alpha}\) of elements of V, all but a finite number of them being 0, such that \(\mathbf{f} - \sum_{\lambda \leqslant \alpha} \mathbf{a}_{\lambda} g_{\lambda} \ll g_{\alpha}\) . One says that \(\sum_{\lambda \leqslant \alpha} \mathbf{a}_{\lambda} g_{\lambda}\) is an asymptotic expansion of \(\mathbf{f}\) to precision \(g_{\alpha}\) , that the \(\mathbf{a}_{\lambda} g_{\lambda} (\lambda \leqslant \alpha)\) are its terms, the \(\mathbf{a}_{\lambda}\) its coefficients, and the function \(\mathbf{r}_{\alpha} = \mathbf{f} - \sum_{\lambda \leqslant \alpha} \mathbf{a}_{\lambda} g_{\lambda}\) the remainder of this expansion.

To express the fact that \(\sum_{\lambda \leqslant \alpha} \mathbf{a}_{\lambda} g_{\lambda}\) is an asymptotic expansion of \(\mathbf{f}\) to precision \(g_{\alpha}\) one most often restricts oneself to writing

\[
\mathbf {f} = \sum_ {\lambda \leqslant \alpha} \mathbf {a} _ {\lambda} g _ {\lambda} + o (g _ {\alpha}) \quad (\text { or } \mathbf {f} = \sum_ {\lambda \leqslant \alpha} \mathbf {a} _ {\lambda} g _ {\lambda} + o _ {k} (g _ {\alpha})
\]

if there are several functions in the proof) following the notation of V, p.~219 and 220.

Of two asymptotic expansions (of two functions, distinct or not) relative to the same scale E, one says that the one with precision of greater index is the more precise.

If \(\sum_{\lambda \leqslant \alpha} \mathbf{a}_{\lambda} g_{\lambda}\) is an asymptotic expansion of \(\mathbf{f}\) to precision \(g_{\alpha}\) , then, for every \(\beta < \alpha\) , \(\sum_{\lambda \leqslant \beta} \mathbf{a}_{\lambda} g_{\lambda}\) is an asymptotic expansion of \(\mathbf{f}\) to precision \(g_{\beta}\) (V, p.~215, prop. 5): one says that it is obtained by reducing the precision of the given expansion \(\sum_{\lambda \leqslant \alpha} \mathbf{a}_{\lambda} g_{\lambda}\) of \(\mathbf{f}\) to \(g_{\beta}\) .

If \(\sum_{\lambda \leqslant \alpha} \mathbf{a}_{\lambda} g_{\lambda}\) and \(\sum_{\lambda \leqslant \alpha} \mathbf{b}_{\lambda} g_{\lambda}\) are asymptotic expansions to the same precision \(g_{\alpha}\) of two functions \(\mathbf{f}_1, \mathbf{f}_2\) , then \(\sum_{\lambda \leqslant \alpha} (\mathbf{a}_{\lambda} + \mathbf{b}_{\lambda}) g_{\lambda}\) is an asymptotic expansion of \(\mathbf{f}_1 + \mathbf{f}_2\) to

precision \(g_{\alpha}\) (V, p.~215, prop. 5); and for every scalar \(c\) , \(\sum_{\lambda \leqslant \alpha} \mathbf{a}_{\lambda}cg_{\lambda}\) is an asymptotic expansion of \(\mathbf{f}_1c\) to precision \(g_{\alpha}\) . It follows that if a function admits an asymptotic expansion to precision \(g_{\alpha}\) then this expansion is unique: it is sufficient to see that the function 0 does not admit an asymptotic expansion with precision \(g_{\alpha}\) having coefficients \(\neq 0\) . For, if \(0 = \sum_{\lambda \leqslant \alpha} \mathbf{a}_{\lambda}g_{\lambda} + \mathbf{r}_{\alpha}\) , and if \(\gamma\) were the least of the indices \(\lambda \leqslant \alpha\) such that \(\mathbf{a}_{\lambda} \neq 0\) , one would have \(\mathbf{a}_{\gamma}g_{\gamma} = -\sum_{\gamma < \lambda \leqslant \alpha} \mathbf{a}_{\lambda}g_{\lambda} - \mathbf{r}_{\alpha} \ll g_{\gamma}\) , which is absurd.

To say that a function \(\mathbf{f}\) admits an asymptotic expansion to precision \(g_{\alpha}\) , all of whose coefficients are zero, is equivalent to saying that \(\mathbf{f} \ll g_{\alpha}\) . If \(\mathbf{f}\) admits an asymptotic expansion \(\sum_{\lambda \leqslant \alpha} \mathbf{a}_{\lambda} g_{\lambda}\) to precision \(g_{\alpha}\) whose coefficients are not all zero, and if \(\gamma\) is the smallest of the indices \(\lambda\) such that \(\mathbf{a}_{\lambda} \neq 0\) , then \(\mathbf{a}_{\gamma} g_{\gamma}\) is the principal part of \(\mathbf{f}\) relative to the scale \(\mathcal{E}\) , for one has \(\mathbf{f} - \mathbf{a}_{\gamma} g_{\gamma} = \sum_{\gamma < \lambda \leqslant \alpha} \mathbf{a}_{\lambda} g_{\lambda} + \mathbf{r}_{\alpha} \ll g_{\gamma}\) ; similarly, if \(\mu \leqslant \alpha\) is an index such that \(\mathbf{a}_{\mu} \neq 0\) , then \(\mathbf{a}_{\mu} g_{\mu}\) is the principal part of \(\mathbf{f} - \sum_{\lambda < \mu} \mathbf{a}_{\lambda} g_{\lambda}\) .

The most important asymptotic expansions in applications are those relative to the scale of the \(x^{-n}\) (resp. of the \(z^{-n}\) ), where n is a positive or negative integer, as x tends to \(+\infty\) or to \(-\infty\) (resp. when the complex number z tends to \(\infty\) ), or relative to the scale of the \((x - c)^{n}\) (resp. \((z - c)^{n}\) ) when the real number x tends to c from the right or left (resp. when the complex number z tends to c). We saw in I, p.~21 that every vector function of a real variable x which is k times differentiable at a point \(c \in R\) admits a Taylor expansion of order k at this point, that is, an asymptotic expansion to precision \((x - c)^{k}\) with respect to the scale of the \((x - c)^{n}\) .

